%% =============================================================================
%% Course Synthesis Companion
%% BCB744 — Quantitative Biology
%% Chapters 1--26
%% =============================================================================
\documentclass[9pt]{article}

\usepackage[a4paper,portrait,top=10mm,bottom=10mm,left=10mm,right=10mm]{geometry}
\usepackage[T1]{fontenc}
\usepackage[utf8]{inputenc}
\usepackage[osf]{ebgaramond}
\usepackage[sfdefault,scaled=0.95]{inter}
\usepackage[scaled=0.88]{FiraMono}
\usepackage{sfmath}
\usepackage[protrusion=true,expansion=true,tracking=true]{microtype}
\usepackage[table]{xcolor}
\usepackage{pagecolor}
\usepackage{longtable}
\usepackage{array}
\usepackage{booktabs}
\usepackage[most]{tcolorbox}
\usepackage{xstring}

\definecolor{MyBackground}{HTML}{FFFDF0}
\definecolor{headbg}{RGB}{28,40,52}
\definecolor{headfg}{RGB}{255,255,255}
\definecolor{groupbg}{RGB}{230,220,200}
\definecolor{diagbg}{RGB}{255,250,230}
\definecolor{flowbg}{RGB}{235,245,240}
\definecolor{rulelight}{RGB}{180,188,198}
\definecolor{mapaccent}{RGB}{0,80,160}

\pagecolor{MyBackground}
\frenchspacing
\pagestyle{empty}
\setlength{\LTpre}{0pt}
\setlength{\LTpost}{0pt}
\setlength{\extrarowheight}{1pt}
\renewcommand{\arraystretch}{1.12}
\newlength{\tableblockwidth}
\setlength{\tableblockwidth}{180mm}
\newlength{\tableoffset}
\setlength{\tableoffset}{10mm}

\newcommand{\Meta}[1]{{\sffamily\small #1}}
\newcommand{\Rc}[1]{{\ttfamily\footnotesize\color{mapaccent}#1}}
\newcommand{\MapCue}[1]{%
  \IfStrEqCase{#1}{%
    {Core}{$\circ$}%
    {G}{$\mu$}%
    {Rk}{$\diamond$}%
    {Asc}{$\leftrightarrow$}%
    {Grp}{$\bullet\!\bullet$}%
    {Bin}{\texttt{0/1}}%
    {Cnt}{\#}%
    {Sm}{$\sim$}%
    {NL}{$f(x)$}%
    {Q}{$Q$}%
    {Pred}{$\hat{y}$}%
    {Work}{$\Delta$}%
  }[\textbf{?}]%
}
\newcommand{\MapRow}[3]{%
  \sffamily\bfseries\color{headbg}#1\\[-1pt]
  {\sffamily\footnotesize\color{mapaccent}\MapCue{#2}\ #2} &
  #3\\[3pt]
}
\newtcolorbox{summarybox}{
  enhanced,
  boxrule=1pt,
  arc=2mm,
  colback=groupbg,
  colframe=headbg,
  colbacktitle=headbg,
  coltitle=headfg,
  title=\textbf{HOW TO USE THIS PAGE},
  fonttitle=\sffamily\bfseries\small,
  left=4mm,right=4mm,top=2mm,bottom=2mm
}

\begin{document}

\noindent\hspace*{\tableoffset}%
\begin{minipage}{\tableblockwidth}
\begin{minipage}[t]{0.56\tableblockwidth}
  \vspace{0pt}
  \raggedright
  {\Huge\sffamily\bfseries\color{headbg} Course Synthesis Companion}\\
  {\large\itshape BCB744 --- Biostatistics Version 3 (3 April 2026)}\\
  {\small Professor A.J. Smit}\\
  {\small Conceptual map of the course storyline. This page complements the procedural cheat sheet and is meant for revision, orientation, and conceptual linking rather than rapid method lookup.}
\end{minipage}
\hfill
\begin{minipage}[t]{0.40\tableblockwidth}
  \vspace{0pt}
  \begin{summarybox}
    \small\sffamily
    Read down the chapter map to recover the course narrative from foundations to model-based workflow.\\
    Use the procedural cheat sheet when you need concrete test choice, assumptions, reporting details, or R functions.
  \end{summarybox}
\end{minipage}
\end{minipage}

\vspace{5pt}
\noindent\hspace*{\tableoffset}%
{\sffamily\footnotesize
\setlength{\tabcolsep}{5pt}
\begin{longtable}{>{\raggedright\arraybackslash}p{34mm} >{\raggedright\arraybackslash}p{136mm}}
\toprule
\rowcolor{headbg}
\color{headfg}\sffamily\bfseries Chapter / Focus & \color{headfg}\sffamily\bfseries Key Learning \\
\endfirsthead
\toprule
\rowcolor{headbg}
\color{headfg}\sffamily\bfseries Chapter / Focus & \color{headfg}\sffamily\bfseries Key Learning \\
\endhead
\midrule
\multicolumn{2}{r}{\sffamily\scriptsize Continued on next page} \\
\endfoot
\bottomrule
\endlastfoot
\midrule
\MapRow{Chapter 1}{Core}{Statistics links biological questions to data, uncertainty, and decisions; population, sample, parameter, and estimate are not interchangeable.}
\MapRow{Chapter 2}{Core}{Descriptive statistics summarise centre, spread, and sample size, but they do not by themselves justify an inferential claim.}
\MapRow{Chapter 3}{Core}{Graphs are part of the analysis, not decoration; choose plots that reveal structure, spread, outliers, and grouping before fitting models.}
\MapRow{Chapter 4}{Core}{Probability distributions and sampling variation explain why estimates fluctuate; standard errors and confidence intervals arise from repeated sampling logic.}
\MapRow{Chapter 5}{Core}{Inference asks whether an observed pattern is large relative to sampling uncertainty; hypotheses, test statistics, \(p\)-values, and intervals answer different parts of that task.}
\MapRow{Chapter 6}{G}{Assumptions attach to the scale and design of the analysis; check normality, equal variance, and transformation needs before committing to Gaussian procedures.}
\MapRow{Chapter 7}{G}{\textit{t}-tests compare means for one-sample, paired, and two-group designs; pairing changes the analysis, and Welch is usually safer than pooled-variance Student's \textit{t}.}
\MapRow{Chapter 8}{G}{ANOVA is the multi-group extension of mean comparison; the omnibus test asks whether any group means differ, and post hoc procedures locate the differences.}
\MapRow{Chapter 9}{Asc}{Correlation quantifies association, not causation; Pearson targets linear association, whereas Spearman and Kendall target monotonic rank association.}
\MapRow{Chapter 10}{Work}{Method selection follows data structure: response type, number of groups, predictor type, pairing, and independence determine the inferential framework.}
\MapRow{Chapter 11}{G}{Regression assumptions are checked on residuals because fitting removes the systematic mean structure; residual plots guide model revision.}
\MapRow{Chapter 12}{G}{Simple linear regression models the conditional mean as a line; slope inference, fitted values, and confidence versus prediction intervals answer different questions.}
\MapRow{Chapter 13}{G}{Polynomial terms allow curvature in the mean function, but coefficients should be interpreted through the fitted curve rather than term by term.}
\MapRow{Chapter 14}{G}{Multiple regression estimates partial effects while adjusting for other predictors; model specification determines which biological questions the coefficients answer.}
\MapRow{Chapter 15}{G}{Interactions test whether the effect of one predictor changes with another; once an interaction is present, main effects cannot be interpreted as universal averages.}
\MapRow{Chapter 16}{G}{Collinearity, confounding, and measurement error destabilise coefficient interpretation; unstable estimates are often a design or variable-definition problem.}
\MapRow{Chapter 17}{Work}{Model checking separates explanation from prediction; compare models with diagnostics, nested tests, information criteria, and out-of-sample error according to the objective.}
\MapRow{Chapter 18}{Grp}{Pseudoreplication is a design error in which non-independent observations are analysed as if they were independent; no later statistical fix rescues bad sampling design.}
\MapRow{Chapter 19}{Grp}{Mixed models represent grouped or repeated-measures data by combining fixed effects with random effects; dependence must be modelled, not ignored.}
\MapRow{Chapter 20}{Bin}{GLMs extend regression beyond Gaussian responses by matching the response distribution and link to counts, proportions, and binary outcomes.}
\MapRow{Chapter 21}{Sm}{GAMs replace rigid functional forms with smooth terms when the response pattern is nonlinear but not well captured by a specific polynomial or parametric curve.}
\MapRow{Chapter 22}{NL}{Nonlinear regression is appropriate when biology implies a specific nonlinear equation; parameter meaning comes from the model form, and starting values often matter.}
\MapRow{Chapter 23}{Q}{Quantile regression models conditional quantiles rather than only the conditional mean, revealing how effects differ across the response distribution.}
\MapRow{Chapter 24}{Pred}{Explanation and prediction are related but distinct goals; the best explanatory model need not be the best predictive model, and vice versa.}
\MapRow{Chapter 25}{Pred}{Regularisation trades some bias for lower variance and better prediction in high-dimensional or collinear settings; ridge shrinks, lasso selects, and elastic net blends both.}
\MapRow{Chapter 26}{Work}{Reproducible workflow is part of statistical practice: script the analysis, record decisions, manage outputs systematically, and keep the inferential path inspectable.}
\end{longtable}
}

\end{document}
